\pdfoutput=1
\documentclass[11pt]{article}
\usepackage[margin=1in]{geometry}
\usepackage{amsmath,amssymb,amsthm,booktabs,array,tabularx,longtable,microtype,adjustbox,xurl}
\usepackage[T1]{fontenc}
\usepackage{lmodern}
\usepackage[round,authoryear]{natbib}
\usepackage[colorlinks=true,linkcolor=blue,citecolor=blue,urlcolor=blue]{hyperref}
\setlength{\emergencystretch}{3em}
\setlength{\parskip}{3pt}
\newcommand{\E}{\mathbb E}
\newcommand{\Prob}{\mathbb P}
\newcommand{\R}{\mathbb R}
\newcommand{\one}{\mathbf 1}
\newcommand{\diag}{\operatorname{diag}}
\newcommand{\TV}{\operatorname{TV}}
\newcommand{\argmax}{\operatorname*{arg\,max}}
\newcommand{\argmin}{\operatorname*{arg\,min}}
\newtheorem{proposition}{Proposition}
\theoremstyle{remark}\newtheorem{example}{Example}
\title{Commercial Value from Correlated Choice Models:\\
Offer Optimization, Event Trading, and AI Service Routing}
\author{Bernd Johannes Wuebben}
\date{September 6, 2026\\[4pt]\small Commercial research and validation plan}
\hypersetup{pdftitle={Commercial Value from Correlated Choice Models},pdfauthor={Bernd Johannes Wuebben}}
\begin{document}
\maketitle
\begin{abstract}
Correlated choice models can support commercial decisions when dependence changes
which action is profitable and the relevant outcomes can be measured. They do
not, by themselves, establish a trading advantage, demand for a product, or an
agent's absolute reliability. We develop three implementable proposals:
contribution-profit optimization for merchants, relative-value trading in event
contracts, and routing among AI services with verification and fallback. A sparse
Gaussian identification result provides a possible way to learn substitution
from controlled comparisons; its experimental assumptions and repeated-sampling
cost are retained. We distinguish exact payoff relationships from speculative
price convergence, and relative performance from successful task completion.
For each proposal we specify the customer or counterparty, decision rule,
required data, incumbent comparison, unit economics, and conditions for stopping.
The provisional first business is a paid merchant experiment on an existing
controlled sales channel. Prediction-market research begins with settlement
semantics and executable costs; broad agent marketplaces and financial guarantees
are deferred. Reproducible calculations quantify discount cannibalization,
sample-size requirements, customer scale, trading capacity, and correlated
fallback failures. These calculations are illustrative assumptions, not measured
commercial results. No strategy backtest, customer pilot, or investment return
is claimed.
\end{abstract}
\vspace{4pt}
\noindent\textbf{Decision in brief.} Obtain control of one economically valuable
decision and its realized outcomes. Sell a measured improvement in that decision.
Use the correlated model only if it improves net value against a simpler policy.

\section{The commercial question}
The relevant question is whether a decision-maker will pay for a measurable
improvement in an existing activity. A merchant already chooses promotions. A
trader already commits capital at quoted prices. An enterprise already pays for
automated work and corrects failures. These activities have observable costs,
incumbent policies, and people responsible for their budgets. They are suitable
starting points for assessing the commercial value of choice models.

The research motivating this paper has two parts. \citet{cotton2021,cotton2026}
study forward probabilities and inversion in stochastic contests: a joint
performance model maps relative means to winning probabilities, and inversion
fits relative means to specified probabilities under a chosen shock law.
\citet{wuebbenident} give an experimental identification result for a stable
Gaussian factor model with idiosyncratic noise. Designed binary interventions
identify the law of performance differences under explicit conditions.
\citet{wuebbencompanion} explain applications and limitations through synthetic
examples. The present paper asks what an operator must add before those results
can support revenue or trading profit.

The commercial contribution is a set of decision specifications and falsifiable
investment cases. It is not a new general theory of pricing or routing. In
particular, the existence of a fast inverse does not establish that the true
dependence structure is known, that a competitor is mispricing an event, or that
a buyer will purchase a new service. A public mathematical result also cannot
be treated as an exclusive commercial asset. A defensible business would need
customer access, permission to collect useful outcomes, reliable execution,
and an improvement large enough to cover acquisition and delivery costs.

\subsection{A provisional ordering}
Table~\ref{tab:ranking} ranks the proposals by the feasibility of demonstrating
economic value. This is a judgment conditional on access, not an estimate of
market size. No merchant, enterprise design partner, trading account, or initial
capital commitment is assumed to exist. The planning case is a founder-led team
with one technical builder and part-time domain support. All financial amounts
are in US dollars; they are not forecasts of local salaries or legal costs.

\begin{table}[htbp]\centering\small
\begin{tabularx}{\textwidth}{p{.19\textwidth}XX}
\toprule
Proposal & First economically meaningful test & Provisional priority \\
\midrule
Merchant offers & Randomized increase in realized contribution profit after
returns, discounts and operating costs & First business if a merchant grants
offer control and sufficient traffic \\
Event-contract trading & Positive net profit under exact settlement rules,
conservative fills and capital charges & Bounded research project; deployment
requires independent evidence of executable advantage \\
AI service routing & Lower total cost at an agreed quality and latency standard
on a buyer's actual workflow & First business instead if enterprise data and a
budget owner are accessible \\
Agent marketplace & Repeat transactions between existing buyers and providers
with a reason to pay an intermediary & Expansion after a routing service has
customers and repeat volume \\
Failure guarantees & Measured loss frequency, severity, dependence and funded
contractual obligations & Defer until failure economics are established \\
\bottomrule
\end{tabularx}
\caption{Access-dependent priorities. These are proposed tests, not completed
experiments or market-size estimates.}\label{tab:ranking}
\end{table}

The default recommendation is to sell a merchant profit experiment on a channel
the merchant already controls. The first product need not depend on AI shopping
traffic. If an enterprise workflow owner is substantially easier to reach,
verified service routing can take its place. Running two full businesses before
either has a paying customer would consume the same limited engineering and
sales resources. A small prediction-market data study may be justified, but it
has its own time and expense limit.

\subsection{Evidence available at the drafting date}
The two local research manuscripts contain synthetic calculations. They contain
no measured merchant response, enterprise task-outcome panel, or live trading
record. The present financial tables are generated from printed planning
assumptions. Public documentation and primary research were checked on
6 September 2026 to establish that the relevant interfaces and competing
approaches exist. Product announcements are evidence of infrastructure, not
evidence of our ability to obtain access or earn a return.

That distinction changes the research sequence. A prototype demonstrating
correct probability arithmetic can be useful, but the first investment decision
should concern data and customer access. There is no basis here for an expected
Sharpe ratio, a company valuation, a claimed addressable market, or a probability
of commercial success.

\section{What the choice model identifies, and what a business still needs}
\subsection{The experimental result}
For a fixed task population or shopping context, consider latent utilities
\begin{equation}\label{eq:rum}
 U_i=\mu_i+v_i^TZ+\sqrt{d_i}\eta_i,\qquad
 Z\sim N(0,I_r),\quad \eta_i\sim N(0,1),
\end{equation}
where $d_i\geq0$ and $Z,\eta_1,\ldots,\eta_n$ are independent.
The covariance is $\Sigma=VV^T+\diag(d)$. An alternative wins a menu when its
utility is greatest. The model assumes that restricting the menu preserves the
joint law of the retained utilities.

If a controlled intervention adds a known $t$ to alternative $i$ without changing
the other utilities or shock law, its binary response against $j$ is
\begin{equation}\label{eq:probit}
 q_{ij}(t)=\Phi\left(\frac{\mu_i-\mu_j+t}{s_{ij}}\right),\qquad
 s_{ij}^2=\Sigma_{ii}+\Sigma_{jj}-2\Sigma_{ij}.
\end{equation}
Two interior population responses at $t_1<t_2$ identify the mean difference and
contrast variance. With $z_\ell=\Phi^{-1}(q_{ij}(t_\ell))$,
\begin{equation}
 s_{ij}=\frac{t_2-t_1}{z_2-z_1},\qquad
 \mu_i-\mu_j=s_{ij}z_1-t_1.
\end{equation}
For complete pair data, define $D_{ij}=s_{ij}^2$, $D_{ii}=0$, and
$P=I-\one\one^T/n$. Then $P\Sigma P=-PDP/2$ determines the Gaussian contrast
law and consequently all menu and ranking probabilities on those alternatives.

The sparse result of \citet{wuebbenident} assumes $r\geq1$,
$n\geq3(r+1)$, positive pairwise contrast variances, and three disjoint groups
of $r+1$ alternatives whose factor loadings are each affinely independent.
Comparing every alternative with those reference alternatives identifies
$P\mu$, $PVV^TP$ and $d$, uniquely within factor dimension at most $r$.
With $k=3(r+1)$ references, the design has
\begin{equation}\label{eq:counts}
 m=kn-\frac{k(k+1)}2
\end{equation}
distinct pairs. The reconstruction takes $O(nr^2+r^3)$ arithmetic and $O(nr)$
storage in factor form. Local statistical stability additionally requires
bounded parameters, positive idiosyncratic variances, informative treatment
responses and uniformly conditioned reference groups. Unknown-rank recovery,
automatic selection of conditioned references, and unrestricted random-pair
recovery with unknown diagonal variances are not established by that theorem.

The number of trials can dominate the arithmetic cost. With $N$ independent
observations at each of two levels per pair, there are $2mN$ binary observations.
Table~\ref{tab:anchorcost} illustrates this distinction for $r=3$. The stated
$N=100$ is a counting assumption; it is not claimed to provide adequate
accuracy for any commercial task. Reusing one sealed score draw at many offsets
does not create independent new task outcomes. If full score vectors are
available, their differences should generally be estimated directly.

\begin{table}[htbp]\centering\small
\begin{tabular}{rrr}
\toprule
Alternatives & Distinct pairs & Binary observations ($N=100$) \\
\midrule
1,000 & 11,922 & 2,384,400 \\
1,000,000 & 11,999,922 & 2,399,984,400 \\
\bottomrule
\end{tabular}

\caption{Observation counts under the minimum three-group design with twelve
references. Two population probabilities per pair require repeated sampling.}
\label{tab:anchorcost}
\end{table}

\subsection{Three gaps between identification and economic value}
First, an observed discount $h$ gives a utility shift $\beta h$ only under an
additive response with a common coefficient $\beta$. A common unknown positive
coefficient permits a shared normalization, but unrestricted customer
heterogeneity, attention effects and perceived-quality changes can invalidate
a pooled Gaussian response. Randomization identifies the effect of the actual
offer policy even when this structural interpretation fails. It does not make
the structural assumptions true.

Second, financial trading does not ordinarily provide this intervention. Buying
an election contract changes inventory and perhaps its price; it does not add
a known amount to the underlying candidate's electoral performance. Calibrating
a contest model from market prices therefore remains model-based inversion.
The sparse experimental theorem does not identify electoral covariance from
order flow.

Third, relative success and absolute success are different outcomes. This
matters directly to the proposed agent exchange and its guarantees.

\begin{proposition}[Rank information does not identify task completion]
Suppose all observations are comparisons of $U_i+t$ with $U_j$, for arbitrary
known offsets and menus. Without an absolute reference or additional outcome
observations, these comparisons do not identify
$\Prob(U_i\geq\tau)$ for a fixed externally defined success threshold $\tau$.
\end{proposition}
\begin{proof}
For any constant $c$, replacing $U$ by $U+c\one$ preserves every comparison
and every ranking, including all additive pair experiments. For a nondegenerate
Gaussian coordinate, $\Prob(U_i+c\geq\tau)$ tends to zero or one as $c$ tends
to minus or plus infinity. The observations are unchanged along this family.
\end{proof}

This is an elementary consequence of the location indeterminacy, rather than
a new identification theorem. A deterministic benchmark with an observed
pass/fail outcome can supply additional information. Its validity must be
established independently: adding a reference label to a leaderboard does not
measure actual remediation cost. Common random shocks are likewise invisible
to differences when added to every score and can change joint absolute
failure probabilities. Correlation of latent ranking scores is not automatically
correlation of operational failures.

\subsection{When greater probability accuracy can pay}
Let $a$ denote an admissible offer or selection policy and $Y$ its outcome.
Suppose its realized contribution $g_a(Y)$ lies in an interval of width at most
$M$, and let $p_a$ and $\widehat p_a$ denote the true and fitted outcome laws.
These distributions include the outside option when no purchase is possible.
For discrete outcomes, our convention is
$\TV(p,\widehat p)=\tfrac12\sum_y|p(y)-\widehat p(y)|$.

\begin{proposition}[A sufficient decision-accuracy bound]\label{prop:value}
If $\TV(p_a,\widehat p_a)\leq\epsilon$ for every action in a fixed finite
action set, choosing an action with the largest fitted expected contribution
loses at most $2M\epsilon$ in true expected contribution relative to the
best action in that set.
\end{proposition}
\begin{proof}
For a function with range width at most $M$, the difference between its
expectations under two distributions is at most $M$ times their total
variation distance. Let $a^*$ maximize true value and $\widehat a$ maximize
fitted value. Insert the fitted values of these two actions between their true
values. The fitted-value difference is nonpositive in the regret calculation;
the two remaining errors are each at most $M\epsilon$.
\end{proof}

The bound is a decision-theoretic diagnostic, not a commercial guarantee.
Actual profit can have refund and loss tails beyond the assumed interval.
Moreover, the local identification paper's fixed-menu bound does not provide
uniform error control for arbitrary prices, contexts or changing menus.
With $M=100$ and $\epsilon=0.01$, the bound permits a loss of \$2 per decision,
which can exceed the improvement being pursued. Reporting lower probability
error is consequently insufficient. The operational test is incremental profit
or total cost on future decisions.

\section{The first business: merchant contribution-profit optimization}
\subsection{Buyer, problem and initial product}
The initial buyer is a head of e-commerce or pricing, with a finance owner who
can approve the profit measurement. The qualifying merchant sells several
substitutable products, can vary offers on an owned channel, has reliable
contribution-cost data, and has enough eligible traffic to run a powered test.
The channel can be a product comparison page, a quote flow, or a merchant-owned
shopping assistant. A store selling only a few high-value items each week is
unlikely to generate the needed experimental evidence quickly.

The proposed product returns one approved offer from a small discrete set,
subject to merchant-specified margin and inventory constraints. It records the
decision, follows the resulting order through returns, and produces a finance-
reconciled treatment-control comparison. Its sales proposition is a measured
increase in retained contribution profit. Conversion alone is a misleading
target when incentives are costly or sales move away from higher-margin items.

The first implementation should use a merchant's existing experimentation and
checkout systems wherever possible. This is already a competitive market.
Intelligems advertises price, discount, shipping and checkout experiments with
cost-of-goods-based metrics \citep{intelligems}. The commercial comparison is
therefore against the merchant's current policy and a competent experiment-led
alternative, not against a fictional operator who discounts at random forever.
The correlated model earns a place only if cross-product substitution improves
decisions beyond those alternatives.

\subsection{The profit objective}
Let $x$ contain information available before the offer: eligible product set,
stock, customer or session attributes that may be used, delivery constraints,
and time. Let $S(x)$ be the observed choice set, with $0$ denoting no purchase
within the prespecified attribution window. For products owned by the merchant,
let $M\subseteq S(x)$ and let $d$ be an admissible discount vector. Define
\begin{equation}\label{eq:merchant}
 V(d,x)=\sum_{i\in M}p_i(d\mid x,S)
             \overline m_i(d,x)-c_{\rm decision}(d,x).
\end{equation}
Here $\overline m_i(d,x)$ is expected retained contribution conditional on buying
$i$ under that policy. It includes collected revenue less discounts, product
cost, payment and channel charges, shipping subsidies, expected returns,
chargebacks and other incremental servicing costs. The decision cost includes
charges incurred even without a purchase. Estimation should use the realized
transaction ledger as the outcome rather than assume these conditional margins
are policy-invariant.

The policy solves $\max_{d\in\mathcal D(x)}V(d,x)$, with an unchanged offer
included in $\mathcal D(x)$. Floors on retained margin, offer ceilings, stock
limits and a restricted action set are explicit business constraints. For a
single product with constant pre-discount contribution $m$ and differentiable
purchase probability $p(d)$, an interior stationary point satisfies
\begin{equation}
 (m-d)p'(d)=p(d).
\end{equation}
This is only a first-order condition. Without additional shape restrictions,
it proves neither uniqueness nor global optimality; the finite approved action
set can instead be evaluated directly.

For multiple products with fixed pre-discount margins and a common utility
response $\beta$ to money, let $L_{ij}=\partial p_i/\partial\mu_j$. Then,
locally and with other discounts held fixed,
\begin{equation}\label{eq:cannibal}
 \frac{\partial V}{\partial d_j}
 =-p_j+\beta\sum_{i\in M}(m_i-d_i)L_{ij}
       -\frac{\partial c_{\rm decision}}{\partial d_j}.
\end{equation}
The sum makes substitution economically relevant. A discount can divert buyers
from a high-margin product owned by the same merchant. For smooth additive
random utility, competing alternatives have nonpositive cross derivatives;
the magnitude and ownership of diverted demand matter. If margins, return
rates or service costs depend on the offer, their derivative terms must also
be included. Equation~\eqref{eq:cannibal} is a restricted explanatory calculation,
not the proposed estimator for an unrestricted transaction ledger.

\begin{table}[htbp]\centering\small
\begin{tabular}{rrrrrr}
\toprule
Discount A & $p_A$ & $p_B$ & No purchase & Conversion & Contribution \\
\midrule
\$0 & 0.040 & 0.060 & 0.900 & 10.0\% & \$8.400 \\
\$5 & 0.065 & 0.052 & 0.883 & 11.7\% & \$8.775 \\
\$10 & 0.090 & 0.030 & 0.880 & 12.0\% & \$7.500 \\
\bottomrule
\end{tabular}

\caption{Illustrative discount cannibalization. A contributes \$60 before its
discount and B contributes \$100. Probabilities are supplied scenario inputs,
not data or a fitted Gaussian demand system. Contribution is per eligible
session, before any common decision charge.}\label{tab:offers}
\end{table}

Table~\ref{tab:offers} gives a simple reason to optimize the merchant's full
contribution. A \$10 discount raises conversion from 10\% to 12\%, yet lowers
expected contribution from \$8.400 to \$7.500 per session. A \$5 discount gives
\$8.775 and is best among the three supplied actions. The example demonstrates
an accounting possibility, not a measured discount response. It establishes
neither Gaussianity nor a likely magnitude of attainable uplift.

\subsection{What can actually be observed}
A merchant's order feed records its own sales, not the losing offers at every
competitor. Nor does a completed checkout reveal every product that a buyer
considered. Accordingly, the first experiment uses a controlled, logged
merchant menu. If competitor offers are missing, a competitor-specific
substitution matrix is not claimed. No purchase is an observed outcome within
the merchant channel; it does not distinguish postponement, abandonment, or
purchase elsewhere without additional data.

The simple winner model concerns one choice. A multi-item basket, subscription
renewal, stockpiling response, or discount that changes future shopping requires
a different outcome model. The first pilot should restrict the setting or
evaluate the complete basket and repeat-purchase contribution directly. It
should not manufacture a single-choice likelihood from a basket transaction.

Prices in a scraped catalog likewise do not establish what the customer saw.
The required menu log records actual exposures, prices, inventory and delivery
terms at the decision time. When that log cannot be obtained, a randomized
offer-policy test may still identify the merchant's overall effect, but the
structural account of who substituted for whom remains limited.

\subsection{AI commerce is a distribution option}
Google's Universal Commerce Protocol describes checkout capabilities and
extensions, including discounts \citep{ucp}. Google's May 2026 announcement
describes an expansion of its Direct Offers pilot \citep{directoffers}. These
interfaces make merchant offer logic a relevant application, but neither
announcement establishes open access to a real-time auction over a buyer's
full consideration set.

OpenAI's September 2025 announcement introduced the Agentic Commerce Protocol
with Stripe \citep{acp}. It should be read as a dated launch announcement.
Shopify's June 2026 description distinguishes channel-specific checkout paths,
including merchant-site checkout through an in-app browser for ChatGPT
\citep{shopifycommerce}. The integration route must therefore be confirmed for
the chosen merchant and channel. A protocol specification does not guarantee
that every checkout customization or randomized offer can be applied there.

The practical first release works on the merchant's existing site or quote
flow. An AI channel is added when its documented interface supports the needed
offer action, exposure log and outcome attribution. This avoids making early
revenue depend on an unconfirmed platform permission. Public growth claims
about AI referrals are not converted into a serviceable market estimate here.

\section{Measuring merchant value and earning the fee}
\subsection{A prospective experiment}
The pilot has four stages: instrumentation, randomized exploration, policy
selection using only accumulated training data, and a separate prospective
comparison of the selected policy with the incumbent. An A/A period first checks
that assignment, exposure, orders and returns reconcile without an artificial
treatment difference. The same customer receives a consistent assignment where
cross-session contamination is plausible. Inventory interference may require
randomization across suitable time blocks or product groups instead; its design
and uncertainty calculation must match the unit of randomization.

The primary outcome is retained contribution per eligible randomized unit,
including units with zero purchases. The denominator is set before treatment.
Order cancellations and returns are included after an agreed maturity window.
Changing the denominator to purchasers would condition on a treatment outcome.
Total contribution across the owned substitute set is measured, so a gain on
one product cannot conceal losses on another.

The policy comparison includes: the existing offer rule, a simple segmented
best-offer policy, a logit or other standard response model where appropriate,
and the correlated model. Models receive the same information available before
the decision. If the correlated model does not increase prospective net profit
after its additional computational and maintenance costs, it is removed from
the commercial product. A profitable simple policy remains a legitimate
business result.

Exploration records the probability $e(a\mid x)$ of every assigned action.
For a fixed candidate policy $\pi$, one possible supplementary offline estimator
is the doubly robust policy-value estimator \citep{dudik2011}:
\begin{equation}\label{eq:dr}
 \widehat V(\pi)=\frac1T\sum_{t=1}^T\left[
 \sum_a\pi(a\mid x_t)\widehat m(x_t,a)
 +\frac{\pi(a_t\mid x_t)}{e(a_t\mid x_t)}
       \{Y_t-\widehat m(x_t,a_t)\}\right].
\end{equation}
Here $Y_t$ is realized contribution and $\widehat m$ is a fitted conditional
mean. This requires consistency of outcomes, action overlap and a defensible
assignment mechanism; nuisance predictions and policy selection are separated
from their evaluation observations. An unoffered action cannot be evaluated
by assigning it an invented propensity. Cluster dependence and adaptive
assignment require corresponding inference. The final customer claim is based
on the prospective policy comparison, not the best-looking value after searching
many policies on the same log.

\subsection{Traffic can be the binding constraint}
For a prespecified comparison of two equally sized independent groups with
common outcome standard deviation $\sigma$, a normal approximation to 80\%
power at a two-sided 5\% level gives
\begin{equation}\label{eq:power}
 n_{\rm arm}\simeq
 \frac{2(z_{0.975}+z_{0.80})^2\sigma^2}{\Delta^2}.
\end{equation}
This is a planning calculation, not an exact test for heavy-tailed profit data.
Table~\ref{tab:power} sets $\sigma=\$26$, approximately the standard deviation
of the unchanged-offer scenario in Table~\ref{tab:offers}. A pilot must replace
this assumption with an estimate from its own historical or A/A data.

\begin{table}[htbp]\centering\small
\begin{tabular}{rrr}
\toprule
Effect/session & Observations/arm & Total observations \\
\midrule
\$0.100 & 1,061,169 & 2,122,338 \\
\$0.250 & 169,787 & 339,574 \\
\$0.375 & 75,461 & 150,922 \\
\bottomrule
\end{tabular}

\caption{Illustrative sample sizes for a fixed two-policy comparison, 80\%
power, two-sided 5\% level and \$26 standard deviation per independent unit.
These do not include policy-learning samples or a multiple-comparison penalty.}
\label{tab:power}
\end{table}

Detecting \$0.10 per session requires over two million independent observations
in this scenario. Detecting the illustrative \$0.375 improvement requires
150,922. If only 10\% of store traffic is eligible, total site traffic is not
the available sample size. Repeated users, seasonal changes, return delays and
inventory dependence further affect feasibility. Small pilots should target a
larger measurable business decision or run longer; they should not describe an
underpowered null result as evidence of no effect.

The selling threshold is economic as well as statistical. Before the comparison,
the parties specify the smallest benefit that covers the customer's switching
and operating costs and the supplier's fee. A lower confidence bound above that
threshold is the proposed rollout condition. A nonsignificant estimate triggers
the previously agreed maximum-duration or maximum-budget decision, rather than
an indefinite sequence of new segment searches.

\subsection{Commercial terms and attribution}
The first contract can charge a fixed \$10,000--\$25,000 instrumentation and
analysis fee, explicitly offered as a negotiating hypothesis. This is a paid
service deliverable and does not promise a positive treatment effect. Payment
willingness is untested. A portion can be credited against a later service
contract if the pilot produces a finance-approved benefit.

For a performance component, let $\Delta\Pi_H$ denote estimated incremental
contribution over an agreed measurement period $H$. A possible fee is
\begin{equation}\label{eq:fee}
 F_H=\tau\max\{0,\underline{\Delta\Pi}_H\},
\end{equation}
where $\underline{\Delta\Pi}_H$ is a prespecified conservative estimate, such as
a valid lower confidence bound, and $\tau$ is the negotiated share. The
measurement schedule and handling of repeated looks are fixed in advance.
Loss carryforwards and periodic reconciliation prevent the supplier from
collecting on favorable noisy periods while discarding unfavorable periods.
An alternative is a fixed recurring fee following a successful experiment;
value-based pricing is not mandatory.

The agreement must define eligibility, assignment, returns, horizon, baseline
changes, taxes, incremental servicing costs and audit rights. For a subscription
product, retaining a randomized control fraction can support continued
measurement, subject to commercial and power considerations. Annual revenue
projections must use the traffic actually receiving the policy, not all store
sales. Training and measurement costs, as well as the foregone gains of the
control group, belong in the total business case.

\subsection{Customer economics and supplier scale}
Let $G$ be annual baseline contribution on eligible traffic to which the new
policy is applied, $u$ its fractional improvement, and $\tau$ the effective fee
share actually earned after measurement rules. A simplified mature account has
\begin{equation}\label{eq:scale}
 R=\tau uG,\qquad
 \Pi_{
\rm supplier}=J(R-c)-F,
\end{equation}
where $J$ is the customer count, $c$ annual direct delivery cost per customer,
and $F$ annual fixed operating cost. This expression excludes one-time pilot
fees and assumes collected fees, stable accounts and no bad debt. In an actual
plan, onboarding, churn, acquisition expense and collection timing are modeled
explicitly. A statistically conservative fee can make realized $R$ lower than
the deterministic scenario.

\begin{table}[htbp]\centering\small
\begin{adjustbox}{max width=\textwidth}\begin{tabular}{rrrrr}
\toprule
Uplift & Annual fee & Direct contribution & Break-even clients & Clients for \$1m \\
\midrule
1\% & \$15,000 & \$3,000 & 200 & 534 \\
3\% & \$45,000 & \$33,000 & 19 & 49 \\
5\% & \$75,000 & \$63,000 & 10 & 26 \\
\bottomrule
\end{tabular}
\end{adjustbox}
\caption{Illustrative mature-account economics: $G=\$10$ million,
$\tau=15\%$, $c=\$12,000$ and $F=\$600,000$ annually. The last column targets
\$1 million annual operating profit before tax, with customer counts rounded up.
The uplifts are assumptions.}\label{tab:merchantecon}
\end{table}

At the 3\% assumption, the customer gains \$300,000 before the fee and retains
\$255,000 afterwards, before its additional internal costs. The supplier earns
\$45,000 revenue and \$33,000 direct contribution per customer. Twenty such
customers produce \$900,000 revenue but only \$60,000 operating profit after
the assumed fixed expense. Forty produce \$720,000 operating profit. Forty-nine
are needed to exceed \$1 million. This is a possible valuable company, but it
requires many substantial eligible accounts or materially better economics.

At \$45,000 per account, \$10 million recurring revenue requires 223 equivalent
accounts and \$100 million requires 2,223. These are arithmetic requirements,
not a count of available buyers. Qualifying merchants may be much scarcer than
the total merchant population because the eligible traffic and contribution
base are restricted. At an assumed 25\% contribution margin, the \$10 million
eligible contribution base requires \$40 million eligible annual sales. If only
10\% of sales are eligible, total sales would need to be \$400 million under
the same margin assumption. The commercial segment must therefore be chosen
from eligible economics rather than a generic list of small online stores.
At a 1\% improvement, direct contribution falls to
\$3,000 per customer, making the same delivery model unattractive.

Customer acquisition can overwhelm the model as well. A hypothetical \$25,000
acquisition and onboarding cash outlay takes about 9.1 months to recover from
\$33,000 annual direct contribution, assuming revenue starts immediately and
the account remains active. A three-month sales and deployment delay extends
cash recovery accordingly. An assumption that direct support costs only
\$12,000 must therefore be tested early. The founder should log implementation
hours and support demands from the first pilot.

\section{Prediction-market trading: a separate investment case}
\subsection{The source of possible profit}
Three different activities must be distinguished. Exact payoff arbitrage buys
a package whose contractual payoff is bounded below by more than its full
cost. Statistical relative value trades a forecast of future prices or outcomes
and can lose. Market making supplies immediacy, earning spreads and possibly
rebates while taking inventory, adverse-selection and operating risk.
They require different evidence and should have separate profit attribution.

There is existing research on prediction-market arbitrage.
\citet{saguillo2025} analyze historical Polymarket inconsistencies and extracted
profits. That establishes prior analysis and competition; it does not establish
current accessible opportunities or the performance of the strategy proposed
here. The appropriate first comparison is a semantics-correct payoff scanner
with conservative execution costs. A Gaussian contest model is unnecessary for
identifying an exact relationship between contractual payoffs.

\subsection{Begin with a small, auditable outcome space}
Choose one venue and one family of contracts referring to the same observable
quantity, measurement source and resolution time. A family of threshold and
interval contracts on a single reported scalar is a suitable research target
if the definitions can be verified. The initial family is selected before
examining its apparent profit. Different sources, revised versus first releases,
time zones, missing observations, cancellation and exceptional resolution
clauses may destroy an apparent equivalence.

Partition all contractually possible outcomes into a finite set $\Omega$.
Let $A_{\omega j}$ be the terminal dollar payoff of one purchasable instrument
$j$ in state $\omega$. YES and NO instruments have their own columns. For a
nonnegative purchase vector $x$ and full acquisition-and-carry cost $K_t(x)$,
the guaranteed terminal surplus of a completely acquired package is
\begin{equation}\label{eq:payoff}
 g_t(x)=\min_{\omega\in\Omega}(Ax)_\omega-K_t(x),\qquad x\geq0.
\end{equation}
Displayed depth, lot sizes, cash and position limits constrain $x$. At fixed
price levels and a suitable linear fee representation, maximizing this surplus
can be formulated as a linear program; discrete lots and nonlinear costs can
require a different optimization. Actual venue costs determine the calculation.
An optimizer does not establish that the state list or settlement definitions
are complete.

For a mutually exclusive exhaustive partition, buying one YES share of each
outcome pays one dollar in every state. For two events satisfying
$A\subseteq B$, buying YES on $B$ and NO on $A$ gives
\begin{equation}\label{eq:implication}
 \one_B+(1-\one_A)\geq1.
\end{equation}
The three possible $(\one_A,\one_B)$ states $(0,0),(0,1),(1,1)$ pay
$1,2,1$, respectively. This uses purchasable long instruments and does not
assume unrestricted short selling. The impossible state $(1,0)$ would pay zero:
if contract wording permits that state, the proposed lower bound fails.

The cash flows must be comparable in maturity and collateral terms. A void
refund may have a different payoff from a binary win. Cross-venue packages
add separate oracle, counterparty, withdrawal and settlement risks and may
require collateral at both venues without netting. They are a later research
stage. An automated language model can suggest related contracts, but the
trading relationship is admitted only after a human-readable payoff map and
versioned rule references have been checked.

\subsection{Fees can consume the entire discrepancy}
Kalshi documents market-data and execution APIs, including a public-data
quick start \citep{kalshiapi,kalshidata}. Its fee document effective
7 July 2026 gives a general taker formula proportional to
$0.07MCP(1-P)$, with contract multiplier $M$, share count $C$ and price $P$,
and lists exceptions and maker charges \citep{kalshifees}. The document's
rounding text and displayed examples should be reconciled against the venue's
current calculator and actual fill ledger before production use. A historical
backtest needs the schedule applicable at the historical trade time.

As an arithmetic illustration, YES and NO asks of \$0.49 each create a gross
\$0.02 package surplus. Applying coefficient $0.07$ to both legs before
rounding gives fees of
\begin{equation}
 2(0.07)(0.49)(0.51)=\$0.034986.
\end{equation}
The package then loses about 1.50 cents before additional costs. This is not an
observed quote or a venue invoice. It demonstrates why a two-cent discrepancy
cannot be judged without its fee and execution context.

Polymarket's documentation likewise distinguishes fee-bearing categories and
fee parameters \citep{polymarketfees}. Neither venue is represented here by a
permanent universal fee assumption. Rebates and temporary rewards are excluded
from the base case and reported separately when actually earned. The model
also includes bid--ask costs, depth, settlement funding, failed hedges and
fees of any intermediary used to access the venue.

\subsection{A completed package and an attempted package are different}
Equation~\eqref{eq:payoff} applies to the acquired package. Independent fills
of several limit orders are not an atomic transaction. Let $g$ be the net gain
if every leg fills, $\ell$ the loss if an incomplete package must be unwound,
and $f$ the probability of that failure. A simplified attempt has expectation
\begin{equation}\label{eq:legging}
 (1-f)g-f\ell.
\end{equation}
With supplied inputs $g=\$0.006$, $f=5\%$ and $\ell=\$0.15$, this is
$-\$0.0018$ per attempt. The numbers are illustrative, but the loss mechanism
is central. Order size is constrained by the weakest leg, measured cancellation
latency, and the ability to terminate or unwind exposure.

A taker simulation consumes the contemporaneously available depth after the
assumed processing and network delay. A maker simulation records queue
position or an explicitly conservative proxy and requires evidence of volume
trading through that position. Touching a displayed price is not sufficient
evidence of a fill. If the data cannot identify queue outcomes, report execution
bounds or collect better data; do not assign favorable fills by default.

Kalshi's documentation separates live and historical endpoints for markets,
trades and candlesticks \citep{kalshihistory}. Their existence is not proof of a
complete historical depth-and-queue record. Current order books document
available bid levels and the YES/NO complement convention \citep{kalshibook}.
A production research log needs feed sequence integrity, receive timestamps,
clock checks and observed outages. Permission to use and retain the selected
feed is a separate deployment prerequisite.

\subsection{A statistical relative-value strategy}
Only after the payoff and execution baseline is functioning should a stochastic
model be evaluated. Let $Y$ be a vector of terminal contract payoffs and
$\widehat q_t=\E_{\widehat P_t}[Y\mid\mathcal F_t]$ a forecast formed from
information available at time $t$. Here $\mathcal F_t$ contains only data that
had arrived by the decision time. A family of mutually exclusive contestants
can use a Gaussian race. Overlapping binary events require a joint state model;
their probabilities cannot be forced to sum to one as if only one could occur.

For a joint finite-state model with probabilities $\pi_t$, marginal values are
$q_t=A^T\pi_t$. Recovering $\pi_t$ from marginal prices may be underidentified.
A regularizer or a Gaussian structure can select one completion but cannot
turn that selection into observed joint information. In the existing companion,
two covariances fit the same first-place probabilities but imply different
ranked-contract values. That is model uncertainty to price and test.

For a long single contract held to resolution, a conservative research rule is
to consider purchase only when
\begin{equation}\label{eq:robustentry}
 \inf_{P\in\mathcal P_t}\E_P[Y_j\mid\mathcal F_t]
   -\operatorname{ask}_t(j,q)-c_t(j,q)>b_t(j,q).
\end{equation}
The ask is the volume-weighted executable acquisition price for quantity $q$,
$c_t$ includes fees and funding not already in the ask, and $b_t$ is a
prespecified reserve for estimation and execution error. The set $\mathcal P_t$
contains models or parameter regions estimated from prior data. Unless it has
validated statistical coverage, the infimum is a stress test, not a confidence
bound. A large model uncertainty set may correctly produce no trade.

For convergence trading, the target changes to the future executable exit
price at a specified horizon. A current residual to a fitted surface supplies
no evidence that the residual will mean-revert. The study must prespecify its
exit horizon, maximum holding time, invalidation rule and news exclusions, then
estimate future net liquidation value on chronological holdouts. Holding to
settlement and exiting on convergence have different capital and loss profiles;
their results are reported separately.

The key baselines are current executable prices, a simple calibrated marginal
forecast, and a low-dimensional event model without the proposed correlated
contest structure. A comparison that feeds the target's own current price into
an inversion and then calls the reconstructed target an independent forecast
is circular. Leave-target-out calibration can prevent this particular defect,
but it still needs genuine future-outcome or future-price validation.

\subsection{Market making needs an order-flow model}
For a specified Gaussian contest, the Jacobian $L=\partial p/\partial\mu$
describes how model probabilities respond to changes in latent means. The
identity $dp=L\,d\mu$ does not determine what latent update follows a trade.
A trade might carry information, hedge an existing position, or express an
uninformative liquidity demand. The update
\begin{equation}
 d\mu=H_t(\hbox{order flow, news, depth, prior state})
\end{equation}
is an additional statistical model to be fitted and tested.

Inventory risk also concerns payoffs rather than the covariance of Gaussian
latent scores. For one exhaustive winner vector with probabilities $p$, the
payoff covariance is $\diag(p)-pp^T$. For overlapping events it is instead
$\E[YY^T]-\E[Y]\E[Y]^T$ under a coherent joint law. These are different from
$\Sigma$ in~\eqref{eq:rum}. Confusing them would produce a faulty hedge.

For illustration, an inventory-adjusted reservation value may subtract the
marginal increase in a chosen convex risk charge $\mathcal R$:
\begin{equation}
 r_j(t,I)=\widehat q_{j,t}
 -\{\mathcal R_t(I+e_j)-\mathcal R_t(I)\}.
\end{equation}
Bid and ask distances from that value must cover conditional adverse selection,
fees, uncertainty and the empirically estimated fill process. A positive spread
is not automatically profitable. Markouts after fills, open inventory losses,
rebates and terminal settlements need separate ledgers. A new operator should
not begin with autonomous quoting across thousands of contracts before those
quantities can be reconciled.

\subsection{Research protocol, capacity and stop conditions}
The first trading research release covers one rule-reviewed family and one
venue. It records metadata, available depth and trades, with a dated fee table.
It evaluates exact packages first and model-based opportunities separately.
The evaluation split is chronological and groups all linked contracts from
an event together. Training windows, normalization, factor selection, contract
relationships and model tuning use only information available before the test
decision. Revised data and late-arriving news use their actual availability
times. Repeated ticks from one event are not independent investment trials.

The performance report includes realized and unrealized net profit, costs by
type, maximum drawdown, average and peak capital tied up, exposure to each
underlying event, failed-package losses, return on committed capital, and
profits by event and date. It reports the number of independent event clusters
and all attempted model variants. A Sharpe ratio, if eventually reported, needs
an appropriate time-series frequency and treatment of autocorrelation; it is
not constructed from the apparent sample size of quotes.

The initial protocol reserves a completely untouched final event/date block.
A result driven by one resolved event, disappearing under observed latency,
or requiring unrecorded fills does not justify deployment. The proposed
deployment condition is positive net value under a prespecified conservative
cost case, positive evidence across independent event clusters, and sufficient
capacity to cover operating expense. Thresholds and the data horizon are fixed
before examining final results; insufficient events produce an inconclusive
study, not a passed strategy.

Let $Q$ count completed normalized strategy units per year, with entry, hedge
and exit legs counted once as a complete economic position. Let $e$ be expected
net dollars per unit after all variable costs, and $F_T$ fixed annual expense.
Then the planning requirement for target profit $\Pi_T$ is
\begin{equation}\label{eq:tradingeconomics}
 Q=\frac{\Pi_T+F_T}{e}.
\end{equation}
The definition prevents multiplying a per-package edge by exchange volume
that counts both sides or each leg several times.

\begin{table}[htbp]\centering\small
\begin{tabular}{rrrr}
\toprule
Net edge & Completed units/year & Addressable units/year & Cash capacity \\
\midrule
\$0.002 & 650m & 65bn & \$3,250,000 \\
\$0.005 & 260m & 26bn & \$1,300,000 \\
\$0.010 & 130m & 13bn & \$650,000 \\
\bottomrule
\end{tabular}

\caption{Illustrative scale needed for \$1 million annual operating profit with
\$300,000 fixed expense. Addressable activity assumes 1\% participation in
comparable units; cash capacity assumes 100 annual turns and \$0.50 tied per
unit. None of these capacity inputs has been measured.}\label{tab:tradeecon}
\end{table}

At half a cent net edge, the requirement is 260 million completed units per
year. The illustrative cash requirement is \$1.3 million at 100 turns, but
\$13 million at ten turns with the same cash per unit. Many multi-leg packages
need more than \$0.50 per normalized unit, and peak inventory can exceed its
average substantially. Capital may be partitioned across venues. A large
headline market volume does not establish that this strategy can trade the
required volume at the assumed edge.

Before live deployment, the actual legal entity, residence, venue eligibility,
contract permissions, data rights and position limits must be checked under
the then-current rules. This paper does not assume account access. The initial
task is a data and execution study; funding an account or placing orders is a
separate investment decision after the study.

\section{AI services: verified routing before a marketplace}
\subsection{Choose a narrow workflow with observable outcomes}
The candidate buyer is an operations owner with a repeated task and a
measurable completion standard. One possible first workflow is invoice
extraction and reconciliation that produces a proposed record for human or
deterministic verification, rather than autonomously approving payments.
The task definition specifies fields, source documents, matching rules,
exceptions and the audit record. A financial-document extraction or software
maintenance workflow with equally observable outcomes could serve instead.

The first product compares a small number of already procured services, routes
eligible tasks, verifies outputs, and escalates failures. It does not need an
open exchange, a token, a universal agent score, or thousands of vendors.
The buyer pays for lower verified total cost under quality and latency
constraints. The budget owner must supply representative task data and accept
the definition of success before model development.

Generic routing already has established alternatives. RouteLLM learns model
routing from preference data \citep{routellm}; Amazon Bedrock documents a
commercial intelligent prompt-routing service \citep{bedrock}. The proposed
distinction would have to be task-specific outcome verification and economically
useful selection of alternative providers, including failure dependence. It is
not the mere existence of a router. These sources are evidence of related
methods and products, not an exhaustive competitor survey or a claim that no
incumbent can supply the proposed workflow.

\subsection{The objective and the role of dependence}
For task features $x$ known before dispatch, define total cost
\begin{equation}\label{eq:routingloss}
 L_i(x)=\hbox{service cost}+\hbox{verification cost}
 +\hbox{delay cost}+\hbox{remediation and downstream loss}.
\end{equation}
A one-provider rule minimizes $\E[L_i\mid x]$ over eligible services subject
to the buyer's quality and latency requirements. Marginal conditional costs
are enough for that decision. Cross-provider covariance need not improve it.
Furthermore, the provider most likely to rank first can have a worse expected
loss because winner frequency ignores the magnitude of losses.

Dependence becomes useful when the action includes multiple providers,
inspection or a fallback. Consider a restricted sequence: run A, detect every
failure perfectly, run B if A fails, then pay a fixed remediation cost $R$ if
both fail. Let $F_A,F_B$ be the failure events on the same task, $c_A,c_B$ the
attempt costs, and $p_A=\Prob(F_A\mid x)$. Under this specific information
structure,
\begin{equation}\label{eq:fallback}
 \E[L_{A\to B}\mid x]
 =c_A+p_Ac_B+\Prob(F_A\cap F_B\mid x)R.
\end{equation}
The B outcome must refer to the actual fallback condition, including any
context passed from A. Reusing B's unconditional pass rate among ordinary
tasks would generally be incorrect.

\begin{table}[htbp]\centering\small
\begin{tabular}{rrr}
\toprule
Joint failure & Failure of B given A failed & Expected cost/task \\
\midrule
1\% & 10\% & \$0.56 \\
9\% & 90\% & \$2.96 \\
\bottomrule
\end{tabular}

\caption{Illustrative fallback economics with marginal failure probabilities
10\% for both providers, $c_A=\$0.20$, $c_B=\$0.60$ and $R=\$30$.
The two joint failure laws satisfy the same marginal rates.}\label{tab:fallback}
\end{table}

The first row of Table~\ref{tab:fallback} corresponds to independent failures.
The second has strongly overlapping failures. Expected total cost rises from
\$0.56 to \$2.96 despite unchanged marginal failure rates. This is the
commercial quantity a complementary-provider model could improve. It is
measured from paired task outcomes and the actual fallback mechanism, not
inferred from marginal leaderboard rankings.

Perfect verification is an intentionally strong simplification. An operational
model needs false acceptances, false rejections, timeouts, correlated outages
and delay to escalation. For example, an undetected error can bypass fallback
and incur downstream loss. A router that saves model fees by increasing such
errors can be much more expensive. Tasks should initially remain within a
workflow where errors are observable before expensive irreversible actions.

\subsection{Evaluation data and commercial terms}
The minimum dataset links task features, timestamped provider and model version,
prompt and tool configuration, charged cost, latency, outputs, verification
results, eventual corrections and reviewer effort. The outcome standard is
agreed before evaluation. A subset of tasks is evaluated by multiple providers
to estimate joint failures. Paired tasks, organization-level shifts and provider
versions are held together in the evaluation split. Updating a model or
verification rule creates a new relevant regime.

Task-conditioned Gaussian scores are a possible comparative model. Their
feature map, score scale, rank and stability must be justified separately; a
linear expression $\mu_i(x)=x^T\theta_i$ does not prove those properties.
If cardinal scores and completion labels are observed, use that information
directly. Sparse winner-only experiments are relevant only when the interface
withholds scores and allows valid offsets. Artificially reducing rich logs to
binary wins would discard information the buyer already paid to obtain.

The prospective comparison is against the cheapest acceptable single provider,
the buyer's incumbent, and a simple validated escalation rule. All costs in
\eqref{eq:routingloss} enter the primary outcome. Predeclared noninferiority
margins for serious errors and latency prevent a low-cost policy from passing
by silently degrading quality. The trial needs sufficient independent tasks
in the risk-relevant strata, not just an aggregate benchmark score.

A paid implementation and evaluation project is a possible first sale. A
recurring fee can then be capped as a share of verified savings. For example,
one million annual tasks, \$0.30 verified net savings per task and a 20\%
share yield \$60,000 annual supplier revenue. These are assumptions, not a
quoted market price. The economics are plausible only if bespoke integrations,
human auditing, security review and ongoing support cost substantially less.
At much lower task volume or savings, a separately sold router may be a feature
of another product rather than a viable stand-alone company.

\subsection{Why underwriting and an exchange come later}
An agent's probability of ranking first cannot price a failure guarantee, as
Proposition~1 shows. Even an observed pass rate needs a loss model. Expected
covered loss is frequency times conditional severity, with additional expense
and capital required for correlated incidents. A minimal premium calculation is
\begin{equation}
 \hbox{premium}\ \geq\ \E[\hbox{covered loss}]
   +\hbox{administration}+\hbox{capital charge}+\hbox{uncertainty allowance}.
\end{equation}
Calling the promise a service credit does not establish that every possible
contract lies outside applicable insurance or other rules. Its legal treatment
depends on actual obligations and jurisdiction and must be resolved before
such a product is offered. An initial bounded remediation service can instead
be costed as an operational activity with explicit limits.

Rare-event evidence is expensive. For $N$ independent, identically distributed
trials with zero failures, an exact one-sided 95\% binomial upper bound is
\begin{equation}\label{eq:zero}
 q_U=1-0.05^{1/N}.
\end{equation}
At least 373 zero-failure trials are needed to put that bound at or below
0.8\%; at least 2,995 are needed for 0.1\%. These are best-case counts at a
fixed sample size, before task heterogeneity, provider changes, multiple claims
or common incidents. They provide no estimate of loss severity and no evidence
about rare catastrophic failures absent from the task sample.

An open marketplace adds provider acquisition, buyer acquisition, identity,
dispute resolution, payments and an incentive to transact through the platform
again. Reputation research recognizes task heterogeneity and evaluation
incentives \citep{agentreputation}; that is not evidence that a new exchange
has demand. The service should expand into marketplace functions only after
existing customers repeatedly request provider discovery or procurement and
will pay for it. A take rate must compensate a continuing service, otherwise
repeat buyers and providers have a reason to contract directly.

\section{Execution as a business plan}
\subsection{Access determines the first product}
The immediate decision is which party will provide an executable pilot. A
qualifying merchant supplies offer control, exposure logs, cost accounting and
sufficient traffic. A qualifying enterprise supplies representative tasks,
verification outcomes, provider access and an accountable budget owner.
Institutional trading relationships may instead make a dealer-inquiry product
more accessible. Cotton's original inquiry application selects dealers while
accounting for inquiry costs \citep{cotton2021}; the existing companion explains
the dependence benefit.

A dealer-inquiry service would require timestamped invitations, responses,
prices, no-quotes, sizes, inventory context and executions, including outcomes
for losing dealers. Public transaction prints alone do not reveal the
counterfactual offers of uninvited dealers. The decision objective would be
execution value after inquiry and information-leakage costs, not an inferred
dealer ranking. This is a conditional alternative if a desk can supply the
required logs and support a controlled evaluation. No such relationship is
assumed in the present plan.

The potential institutional economics can be expressed without inventing an
execution result. If a desk routes annual notional $V$ and a controlled study
finds improvement $b$ basis points in its all-in execution cost, measured in
dollars per dollar of notional, annual savings are $Vb/10{,}000$. With assumed
$V=\$50$ billion and $b=0.2$, savings are \$1 million; a 15\% fee would be
\$150,000 before supplier costs. This illustrates why a few large workflows
can support a service. The improvement, routed volume and willingness to pay
remain unmeasured, and the comparison must include no-quotes, lost fills,
latency and the cost of revealing trading interest. A change in dealer
selection without a credible counterfactual execution benchmark is insufficient.

The first two weeks should therefore establish whether a party will supply
the essential data and pay for a defined evaluation. Expressions of interest
without access or a budget owner are weak evidence. Building a universal
platform before resolving this question is not justified by the mathematics.

\subsection{A six-month capital allocation}
Table~\ref{tab:budget} supplies one explicit budget for the main business
experiment. It includes a founder draw and external technical capacity; it is
a planning envelope rather than a market salary survey. The founder may be
able to reduce cash costs, but unpaid labor still has an opportunity cost.
The budget excludes trading capital, insurance reserves and customer-funded
discounts. Merchant incentives and support costs are nevertheless included in
the customer's measured economic outcome.

\begin{table}[htbp]\centering\small
\begin{tabular}{lr}
\toprule
Six-month expense assumption & Amount \\
\midrule
Founder draw & \$60,000 \\
Engineering and data & \$80,000 \\
Experimentation and compute & \$15,000 \\
Customer integration and domain support & \$20,000 \\
Contracts and security & \$15,000 \\
Contingency & \$10,000 \\
Total & \$200,000 \\
\bottomrule
\end{tabular}

\caption{Illustrative six-month cash budget for one main business pilot and
initial product. No customer revenue is assumed to finance it.}
\label{tab:budget}
\end{table}

The \$200,000 envelope is released by milestones. At most \$20,000 is used
to establish access, data feasibility and a written pilot scope. The next
\$40,000 funds instrumentation, baseline measurement and a functioning
experiment. A further \$60,000 funds the prospective evaluation after those
steps succeed. The remaining \$80,000 funds productization and a second
customer only if measured value and delivery economics justify it. These
tranches sum to the budget; the release conditions matter more than the
calendar. Failure to satisfy them preserves the unspent balance.

An optional prediction-market study is capped separately at \$20,000 and
four to six weeks of research effort. It must not consume the main customer's
implementation capacity. The amount is research expense, with no trading
balance included. A capital allocation for actual trading is specified only
after fill, capacity and loss evidence exist. Launching both a trading desk
and a software company within the same \$200,000 assumption is not proposed.

\subsection{Milestones and decisions}
\begin{longtable}{p{.13\textwidth}p{.39\textwidth}p{.39\textwidth}}
\toprule
Period & Required deliverable & Investment decision \\
\midrule\endhead
Weeks 1--2 & A proposed list of 15 relevant buyer conversations; at least one
budget owner supplies a representative data extract and agrees a priced pilot
scope, measurement definition and access path & Without qualified access,
change the customer segment or pause before building the optimizer. The
conversation target is a plan, not a completed sales funnel. \\
Weeks 3--6 & Cost ledger, exposure or task log, A/A or equivalent pipeline
check, simple baseline, sample-size estimate, and maximum trial budget &
Proceed only if the primary outcome is measurable and the attainable sample
can detect an economically worthwhile effect. \\
Weeks 7--12 & Randomized exploration and a locked prospective policy
comparison; models and parameters dated before evaluation & A validated
simple-policy benefit can justify continuation. Without model benefit, omit
the correlated model. Without measurable business benefit, stop or redesign
within the remaining budget. \\
Months 4--6 & Mature outcomes, agreed profit or cost reconciliation, invoice
and collection, second deployment, recorded integration/support hours &
Proceed to recurring product investment only if a customer pays, the benefit
survives the full cost ledger, and a second account can be served economically. \\
After month 6 & Repeat value after changing conditions, renewal evidence,
and a bottom-up account list with acquisition economics & Expand distribution
and product scope only as evidence supports them; marketplace and guarantee
functions remain separate decisions. \\
\bottomrule
\end{longtable}

Calendar dates do not override the outcome horizon. Returns or sufficient
independent events may take longer than twelve weeks. A documented extension
must have a remaining information objective, cost ceiling and stopping date.
The relevant milestone is information acquired per dollar spent, rather than
the number of model features completed.

\subsection{Customer acquisition and defensibility}
The first sale should be direct and specific: an evaluation of a known
promotion or service-routing decision with a capped scope and a shared outcome
definition. Likely introductions are through a relevant commerce agency,
operations integrator, or existing institutional relationship. These are
proposed channels; no partnership is claimed. The founder records the budget
owner, present solution, data access, delivery effort, procurement steps and
reason for rejection for every serious conversation.

After two implementations, the product should standardize the smallest
recurring workflow: data connector, decision constraints, policy execution,
measurement and reconciliation. If each new customer requires a different
research project, pricing must support a consultancy or the target segment must
narrow. A consultancy can produce profit, but it has different capacity and
valuation economics from repeatable software revenue. Labeling custom labor
as recurring software does not resolve that limitation.

Defensibility would come from permissioned outcome histories, task or product
taxonomy, reliable integrations, accumulated experimental knowledge and a
measurement process the customer's finance team trusts. The theorem itself is
public. Data do not automatically compound across customers: permissions,
confidentiality and distribution differences may require tenant-specific
models. A lower-cost platform can integrate a similar optimizer, so the
customer's recurring reason to buy must be measured and revisited.

\section{What would justify a larger commitment}
The research program has a clear hierarchy of evidence. Mathematical
identification establishes what specified experiments can learn under a model.
Held-out prediction assesses that model on new outcomes. A prospective policy
experiment measures economic improvement. A paid and renewed contract measures
whether a customer values the service after delivery costs. Trading requires
the analogous progression through executable prices, fills, risk, capacity and
realized capital use. None of these steps substitutes for the next.

The merchant proposal is the provisional first choice because its decision
can be randomized and its value reconciled on a channel the customer controls.
The opportunity remains conditional on traffic, data access, attainable effect
size and acquisition economics. AI service routing can be preferable where
enterprise access is stronger and absolute outcomes are verifiable.
Prediction-market research has accessible technical interfaces, but no trading
advantage follows from coherence alone. Its first commercial question is
whether enough net executable edge remains after realistic costs.

Three research results would materially strengthen the case: prospective
evidence that correlated substitution improves profit beyond standard offer
experimentation; measured savings from complementary fallback providers after
verification errors; and event-contract returns that survive settlement review,
chronological holdouts and conservative execution. Until those observations
exist, a broad exchange or a balance-sheet guarantee would add obligations
faster than information. The next commitment should purchase the missing
evidence in a narrow workflow with a paying counterparty.

\appendix
\section{Minimum data contracts}
The following schemas describe requirements for a future pilot. No customer
or exchange dataset with these fields is supplied by this paper.

\paragraph{Merchant decisions.}
Record a unique decision identifier; pseudonymous assignment unit; UTC
assignment and exposure times; eligible menu and offer variants; all displayed
prices, stock and delivery terms; predecision features; policy version and
assignment probabilities; selected action; order identifiers; and a complete
outcome window, including no order. Join order-level collected revenue,
discounts, cost of goods, shipping, fees, cancellations, refunds and return
processing costs. Store both event time and availability time where they differ.
Log exposures before observing conversion. Retain the allocation history when
a customer returns. Keep model-development extracts separate from the final
prospective measurement sample.

\paragraph{Event-contract decisions.}
Record venue and contract identifiers; exact rule text or retained permitted
reference and its version; event-family membership; exhaustive outcome-state
mapping; resolution time, source and exceptional-state treatment; market
status; fee schedule and effective date; prices and quantities at every used
depth level; exchange and local receive timestamps; feed sequence; trades;
simulated or actual order lifecycle; cancellation delays; partial fills; paid
fees; collateral; and final settlement cash flows. Simulation, actual order and
settlement records must have distinct fields. A final report must not silently
replace unsuccessful simulations with later better prices.

\paragraph{AI tasks.}
Record task identifier and organization; arrival time; permitted predecision
features; provider, model and workflow version; supplied tools; routing
probability; invoice cost; start and end time; timeout; output; verifier
decision; blinded adjudication where available; later correction; human time;
and downstream loss or a separately labelled unresolved outcome. A paired
evaluation must record whether B saw A's output or diagnosis. This changes the
meaning of the joint failure law.

\section{Questions for the first buyer meetings}
The proposed interview asks for concrete operating evidence. Which promotion,
inquiry or task-routing decision is made repeatedly? Who owns its budget and
can approve a controlled change? What is the current rule and which product
or internal team already supports it? Which outcomes and costs are recorded,
including failures, no-purchases and returns? How much eligible volume occurs
in a month? Can the parties reconcile a representative outcome sample to the
accounting system? What decisions cannot be randomized, and why? What fee and
implementation effort would be acceptable if the agreed benefit were measured?

A promising meeting ends with a named data owner, a feasible experiment,
a costed implementation scope and a paid deliverable. The absence of any one
does not prove that the entire opportunity is false, but it identifies the
next uncertainty to resolve. A proposed conversion funnel of fifteen
conversations to one pilot is a work target, not a forecasted sales conversion
rate. No outreach has been undertaken for this paper.

\section{Commercial hypotheses and validation requirements}
The commercial proposals concern model-based event trading, merchant offers,
agent selection and guarantees. Their investment cases depend on the
following distinctions and validation requirements.
\begin{enumerate}
\item Structural relationships among events are divided into exact payoff
constraints and statistical dependence. Only the former give contractual
lower bounds, and those bounds apply after acquiring the package.
\item The utility Jacobian is used as a sensitivity conditional on a chosen
model. The mapping from order flow to latent information is a separate
unproved and unvalidated component.
\item Discounts are actual interventions on offers. Their interpretation as
known common-scale utility shifts requires additional assumptions, while a
randomized policy comparison can assess profit without those assumptions.
\item Agent rankings are supplemented by absolute completion, cost, latency
and severity data. A score-offset experiment does not identify those targets
merely by recovering relative Gaussian performance.
\item Linear-in-catalog pair counts are reported together with repeated trials,
conditioning requirements and the cost of new task regimes. There is no
claim that a million-agent service can be rated cheaply from a few isolated
comparisons per agent.
\item Company scale is expressed through collected account revenue, direct
costs, acquisition effort and customer counts. Large market narratives and
speculative valuation labels are replaced by conditions a pilot can test.
\end{enumerate}

\section{Reproducibility and source boundaries}
All numerical tables are generated by the accompanying economic calculation
script from a single assumption file. The program checks probability
normalization, the discount accounting, valid joint-failure probabilities,
the sample-size formulas, exact implication payoffs and observation counts.
Inputs and script hashes are stored with the resulting tables. These are
deterministic calculations, not simulations of observed trading or customer
behavior. Changing the assumptions changes the projections; no input is
estimated from commercial data.

From the parent project directory, reproduce the calculations and document with
\begin{verbatim}
.venv/bin/python commercial/src/economics.py
latexmk -pdf -interaction=nonstopmode -halt-on-error \
  -cd commercial/main.tex
\end{verbatim}
The accompanying source register identifies each external claim, its primary
source and the date consulted. Live interface and fee descriptions may change.
Historical results cited from other authors have not been independently
replicated here. No claim of exhaustive novelty, commercial diligence, legal
clearance or empirical validation is made. The original mathematics
and application manuscripts are separate works and retain their own proof and
review boundaries.

\bibliographystyle{plainnat}
\bibliography{references}
\end{document}
